% Parte comum a todos os briefings: contexto do projeto, decisões e porquês.
% Números só por macros (\n...), nunca digitados.

\section{O projeto em um minuto}
\caixa{\textbf{Pergunta:} quais características socioeconômicas, demográficas e geográficas estão associadas à taxa de homicídios nos municípios brasileiros?}

\begin{itemize}
\item \textbf{Tema:} ODS 16 (paz e justiça), indicador 16.1.1: homicídios por 100 mil habitantes.
\item \textbf{Unidade de análise:} o município. São \nMunBR{} municípios, com \nHomBR{} homicídios em 2021--2023 (\nTaxaBR{} por 100 mil hab./ano).
\item \textbf{Resposta:} taxa de homicídios (agressão + intervenção legal, por residência da vítima), média de 2021 a 2023, do SIM/DATASUS.
\item \textbf{Explicativas:} Censo 2022 do IBGE (renda, desigualdade, trabalho, frequência escolar, urbanização, cor ou raça, idade, migração), INEP 2022 (abandono e distorção idade-série no ensino médio), mais estado, Amazônia Legal e faixa de fronteira.
\item \textbf{Modelo:} regressão linear múltipla para ln(taxa), por mínimos quadrados, nos \nMunEsc{} municípios com \nPopMin{} mil habitantes ou mais.
\end{itemize}

\section{As decisões e o porquê de cada uma}
\small
{\renewcommand{\arraystretch}{1.5}\rowcolors{2}{white}{fundo}
\begin{tabularx}{\linewidth}{@{}>{\raggedright\arraybackslash}p{3.6cm}>{\raggedright\arraybackslash}X>{\raggedright\arraybackslash}p{2.6cm}@{}}
\toprule
\textbf{Decisão} & \textbf{Por quê} & \textbf{Apostila}\\
\midrule
Escopo: municípios com \nPopMin{} mil hab. ou mais & Nos pequenos o homicídio é raro: abaixo de 5 mil hab. a mediana é \nPorteHomMenor{} caso em 3 anos e a taxa oscila $\pm$\nPorteOscMenor\% só pelo acaso; \nZerosBR{} municípios têm taxa zero. O escopo fica com \nPctPopEsc\% da população e nenhum zero & Cap. 1: escopo do modelo\\
Resposta em ln(taxa) & Box-Cox estimou $\lambda$ = \nLambda{} (IC de \nLambdaInf{} a \nLambdaSup); usamos $\lambda = 0$, simples de interpretar. A normalidade melhora ($W$ de \nSWy{} para \nSWln) & Cap. 4: Box-Cox\\
Só a renda (sem \% pobres e alfabetização) & As três medem quase o mesmo: VIF da renda \nVIFrenda{} e sinais contraditórios & Cap. 6: multicolinearidade e VIF\\
Renda e renda², com a renda centralizada & A taxa sobe até a renda média e depois cai (U invertido); centralizar evita a colinearidade entre $X$ e $X^2$ & Cap. 7: modelos polinomiais\\
Estado (27 UFs) no lugar de região & Estados diferem muito (polícia, facções). Com estado: $R^2_a$ de \nRdoisAregiao{} para \nRdoisA, $F$ parcial \nFuf{} e Durbin-Watson de \nDWregiao{} para \nDW{} (vizinhos deixam de ter resíduos parecidos) & Cap. 6 e 7: $F$ parcial, variáveis categóricas\\
Bloco de educação e migração & Frequência escolar 15--17, abandono, distorção idade-série e migração recente; a frequência foi a mais forte entre 26 candidatas testadas. $F$ parcial \nFblocoEd{} e $R^2_a$ de \nRdoisAsemEd{} para \nRdoisAcomEd & Cap. 6: $F$ parcial\\
Amazônia Legal fora do modelo & Com o estado no modelo, ela só varia dentro do Maranhão: fica redundante & ---\\
Variáveis frágeis mantidas por enquanto & \% homens 15--29, \% intenção indeterminada e desocupação saem (ou não) na seleção de variáveis & Cap. 8: seleção\\
Outliers mantidos & A validação dos dados não achou erro; a apostila manda remover só com evidência de erro & Cap. 4: outliers\\
Nada fora da apostila & Regra do grupo: técnicas de fora (pesos, modelos espaciais, de contagem) só com decisão do grupo & ---\\
\bottomrule
\end{tabularx}}
\normalsize

\section{O que o modelo mostra}
\begin{itemize}
\item O modelo explica \nRdoisA{} da variação do ln(taxa) ($R^2_a$); $F$ global = \nFglobal{} com \nGlUm{} e \nGlDois{} gl.
\item \textbf{Renda em U invertido:} a taxa esperada é máxima com renda per capita de cerca de R\$ \nPico{} e cai nos mais pobres e nos mais ricos.
\item \textbf{Mais homicídios} com mais \% preta ou parda, mais urbanização, mais desigualdade, mais abandono e atraso escolar e mais migração recente; \textbf{menos} com mais jovens de 15--17 na escola (efeitos em \% na tabela abaixo).
\item \textbf{Não significativas} no modelo atual: fronteira, densidade e \% homens 15--29.
\end{itemize}

\begin{center}\small
\renewcommand{\arraystretch}{1.1}
\begin{tabular}{@{}cllrcr@{}}
\toprule
\textbf{Hip.} & \textbf{Variável} & \textbf{Mudança} & \textbf{Taxa} & \textbf{IC 95\%} & \textbf{$p$}\\
\midrule
 & Renda & --- & U invertido & --- & $<$0,001\\
 & \% preta ou parda & +10 pontos & $+$19,6\% & [$+$15,8; $+$23,6] & $<$0,001\\
 & \% urbana & +10 pontos & $+$9,4\% & [$+$6,5; $+$12,3] & $<$0,001\\
H2 & Distorção idade-série (EM) & +10 pontos & $+$6,1\% & [$+$1,5; $+$10,9] & 0,009\\
H2 & Frequência escolar 15–17 & +5 pontos & $-$5,7\% & [$-$8,3; $-$3,1] & $<$0,001\\
H2 & Abandono no ensino médio & +5 pontos & $+$4,3\% & [$+$0,6; $+$8,2] & 0,022\\
H1 & Desigualdade & +0,1 & $+$3,2\% & [$+$0,7; $+$5,9] & 0,013\\
H3 & Migração recente & +5 pontos & $+$3,0\% & [$+$0,6; $+$5,5] & 0,014\\
 & Desocupação & +1 ponto & $-$1,4\% & [$-$2,5; $-$0,2] & 0,024\\
 & \% intenção indeterminada & +1 ponto & $-$1,2\% & [$-$1,5; $-$0,8] & $<$0,001\\
\bottomrule
\end{tabular}

\end{center}
{\footnotesize Efeito = variação da taxa para a mudança indicada, com as demais variáveis constantes (100 × $b$ com a resposta em log).}

\section{O que ainda não está resolvido}
\begin{itemize}
\item \textbf{Variância não constante:} Breusch-Pagan = \nBP{} (rejeita). Os coeficientes continuam certos na média, mas os valores p ficam aproximados.
\item \textbf{Seleção de variáveis} ($C_p$, AIC, BIC, PRESS) ainda não foi feita: é da entrega final.
\item Testadas e deixadas de fora (não ajudam): garimpo, desmatamento, Bolsa Família, mães solo e gasto municipal em segurança e assistência (SICONFI 2022).
\end{itemize}

\section{Quem fala o quê}
\small
\begin{tabularx}{\linewidth}{@{}>{\raggedright\arraybackslash}p{2.6cm}>{\raggedright\arraybackslash}X@{}}
\toprule
\textbf{Quem} & \textbf{Slides, na ordem}\\
\midrule
\textbf{Kauan} & Capa; Roteiro; 1 · Questão de pesquisa; Limitações dos dados; Escopo: municípios com \nPopMin{} mil hab. ou mais\\
\textbf{Fabio} & 2 · Base de dados; 3 · Variáveis\\
\textbf{Douglas} & 4 · A taxa é assimétrica; o log corrige; 4 · A violência se concentra no Norte e Nordeste; 4 · Nenhum fator sozinho explica a taxa\\
\textbf{Luiz Felipe} & 4 · A renda tem efeito em U invertido; 5 · Box-Cox indica o logaritmo da taxa; 5 · Construção do modelo em blocos\\
\textbf{Willian} & 5 · Modelo preliminar: efeito de cada variável na taxa; 5 · Efeito da renda previsto pelo modelo; 5 · Diagnóstico dos resíduos; Próximos passos; Obrigado\\
\bottomrule
\end{tabularx}
\normalsize


\section{Glossário rápido}
\small
\begin{tabularx}{\linewidth}{@{}>{\bfseries}p{3.4cm}X@{}}
Escopo do modelo & Faixa de dados em que o modelo vale. Fora dela, as conclusões não se aplicam.\\
Box-Cox & Procedimento que escolhe a potência $\lambda$ da resposta que deixa o modelo mais perto das hipóteses; $\lambda = 0$ é o log.\\
$R^2_a$ (ajustado) & Fração da variação explicada, com penalidade por variável a mais.\\
Teste $F$ parcial & Compara o modelo com e sem um bloco de variáveis: diz se o bloco ajuda.\\
VIF & Mede o quanto uma variável repete as outras (multicolinearidade). Acima de 10 é grave.\\
Durbin-Watson & Testa se os resíduos de vizinhos se parecem. Perto de 2 = independentes.\\
Breusch-Pagan & Testa se a variância dos resíduos é constante.\\
Shapiro-Wilk ($W$) & Testa a normalidade dos resíduos; perto de 1 = mais normal.\\
Outlier & Resíduo semi-studentizado com $|e^*| \geq 4$.\\
\end{tabularx}
\normalsize
